% dynamics/dynamic_balance.tex

\chapter{Dynamic Balance and Operational Admissibility}

\label{chap:dynamic-balance}

% ============================================================
\section{Motivation}
% ============================================================

\begin{remark}
Railway operation involves continuous interaction between:
\begin{itemize}
\item curvature,
\item cant,
\item vehicle response,
\item and operational velocity.
\end{itemize}
\end{remark}

\begin{remark}
Classical railway engineering often expresses this interaction through
equilibrium relations between curvature and cant.
\end{remark}

\begin{remark}
Within AIM, however, such relations are interpreted not as isolated
engineering formulas, but as runtime-state coupling laws.
\end{remark}

\begin{principle}[Dynamic-balance principle]
\label{princ:dynamic-balance}
Within AIM, operational balance emerges from runtime coupling between:
\begin{itemize}
\item curvature evolution,
\item cant evolution,
\item vehicle-state response,
\item realization,
\item and operational context.
\end{itemize}
\end{principle}

% ============================================================
\section{Equilibrium Relation}
% ============================================================

\begin{definition}[Equilibrium relation]
\label{def:equilibrium-relation}
A dynamically balanced runtime state approximately satisfies:
\[
v^2\kappa(s)
=
g\sin\phi(s).
\]
\end{definition}

% ------------------------------------------------------------

\begin{remark}
This relation couples:
\begin{itemize}
\item curvature-induced acceleration,
\item and cant-induced compensation.
\end{itemize}
\end{remark}

% ------------------------------------------------------------

\begin{remark}
Within AIM, the equilibrium relation is interpreted as a runtime-state
coupling law rather than merely as a static engineering equation.
\end{remark}

% ============================================================
\section{Operational Imbalance}
% ============================================================

\begin{definition}[Operational imbalance]
\label{def:operational-imbalance}
Operational imbalance is the deviation between:
\begin{itemize}
\item curvature-induced acceleration,
\item and cant-induced equilibrium compensation.
\end{itemize}
\end{definition}

% ------------------------------------------------------------

\begin{remark}
The canonical imbalance quantity is:
\[
a_{\mathrm{eff}}(s)
=
v^2\kappa(s)
-
g\sin\phi(s).
\]
\end{remark}

% ------------------------------------------------------------

\begin{remark}
Operational imbalance is unavoidable in practical railway operation and
must therefore be interpreted dynamically rather than eliminated
completely.
\end{remark}

% ============================================================
\section{Cant Deficiency}
% ============================================================

\begin{definition}[Cant deficiency]
\label{def:dynamic-cant-deficiency}
Cant deficiency denotes insufficient cant relative to the operational
equilibrium state.
\end{definition}

% ------------------------------------------------------------

\begin{remark}
Cant deficiency corresponds operationally to:
\[
v^2\kappa(s)
>
g\sin\phi(s).
\]
\end{remark}

% ------------------------------------------------------------

\begin{remark}
In this situation, uncompensated lateral acceleration remains within the
vehicle state.
\end{remark}

% ------------------------------------------------------------

\begin{remark}
Cant deficiency influences:
\begin{itemize}
\item passenger comfort,
\item wheel unloading,
\item lateral force distribution,
\item and operational admissibility.
\end{itemize}
\end{remark}

% ============================================================
\section{Excess Cant}
% ============================================================

\begin{definition}[Excess cant]
\label{def:excess-cant}
Excess cant denotes cant exceeding the equilibrium requirement for the
current operational state.
\end{definition}

% ------------------------------------------------------------

\begin{remark}
Operationally:
\[
v^2\kappa(s)
<
g\sin\phi(s).
\]
\end{remark}

% ------------------------------------------------------------

\begin{remark}
Excess cant may lead to:
\begin{itemize}
\item inward acceleration effects,
\item reduced passenger comfort,
\item operational limitations at low speed,
\item and wheel unloading effects.
\end{itemize}
\end{remark}

% ============================================================
\section{Equilibrium States}
% ============================================================

\begin{definition}[Dynamic equilibrium state]
\label{def:dynamic-equilibrium-state}
A dynamic equilibrium state is a runtime state satisfying operational
balance within admissible tolerances.
\end{definition}

% ------------------------------------------------------------

\begin{remark}
Within practical railway operation, equilibrium is not exact but exists
within operational ranges.
\end{remark}

% ------------------------------------------------------------

\begin{remark}
Thus, equilibrium should be interpreted as a band of admissible runtime
states rather than as a single exact condition.
\end{remark}

% ============================================================
\section{Equilibrium Bands}
% ============================================================

\begin{definition}[Equilibrium band]
\label{def:equilibrium-band}
An equilibrium band is a range of admissible operational imbalance
states satisfying:
\[
|a_{\mathrm{eff}}(s)|
\le
a_{\max}.
\]
\end{definition}

% ------------------------------------------------------------

\begin{remark}
Equilibrium bands depend on:
\begin{itemize}
\item vehicle type,
\item operational velocity,
\item standards,
\item realization quality,
\item and comfort requirements.
\end{itemize}
\end{remark}

% ------------------------------------------------------------

\begin{remark}
Within AIM, operational equilibrium is therefore interpreted as a
runtime admissibility region.
\end{remark}

% ============================================================
\section{Dynamic Admissibility}
% ============================================================

\begin{definition}[Dynamic admissibility]
\label{def:dynamic-admissibility}
A runtime state is dynamically admissible if operational quantities
remain within admissible bounds.
\end{definition}

% ------------------------------------------------------------

\begin{remark}
Relevant admissibility quantities include:
\begin{itemize}
\item effective acceleration,
\item jerk,
\item roll evolution,
\item cant deficiency,
\item wheel unloading,
\item and operational smoothness.
\end{itemize}
\end{remark}

% ------------------------------------------------------------

\begin{remark}
Dynamic admissibility therefore depends on coupled runtime-state
evolution rather than on geometry alone.
\end{remark}

% ============================================================
\section{Operational Envelopes}
% ============================================================

\begin{definition}[Operational envelope]
\label{def:operational-envelope}
An operational envelope is the admissible region of runtime operational
states.
\end{definition}

% ------------------------------------------------------------

\begin{remark}
Operational envelopes may constrain:
\begin{itemize}
\item acceleration,
\item jerk,
\item roll behaviour,
\item cant deficiency,
\item speed,
\item and realization tolerances.
\end{itemize}
\end{remark}

% ------------------------------------------------------------

\begin{remark}
Within AIM, operational envelopes define admissible runtime-state
regions rather than purely geometric limits.
\end{remark}

% ============================================================
\section{Dynamic Coupling}
% ============================================================

\begin{remark}
Dynamic balance is fundamentally a coupled runtime phenomenon.
\end{remark}

% ------------------------------------------------------------

\begin{remark}
In particular:
\begin{itemize}
\item curvature influences acceleration,
\item cant influences roll response,
\item realization influences smoothness,
\item and operational speed influences admissibility.
\end{itemize}
\end{remark}

% ------------------------------------------------------------

\begin{remark}
Thus:
\[
\kappa(s),
\qquad
\phi(s),
\qquad
v(s)
\]
cannot be interpreted independently.
\end{remark}

% ============================================================
\section{Balance Trajectories}
% ============================================================

\begin{definition}[Balance trajectory]
\label{def:dynamic-balance-trajectory}
A balance trajectory is the runtime evolution of operational balance
states along the railway alignment.
\end{definition}

% ------------------------------------------------------------

\begin{remark}
Balance trajectories are generated jointly by:
\begin{itemize}
\item curvature evolution,
\item cant evolution,
\item realization,
\item and vehicle-state response.
\end{itemize}
\end{remark}

% ------------------------------------------------------------

\begin{remark}
Operational smoothness therefore depends on:
\[
a_{\mathrm{eff}}(s),
\qquad
j(s),
\qquad
\phi'(s),
\]
rather than solely on geometric continuity.
\end{remark}

% ============================================================
\section{Runtime-State Coupling}
% ============================================================

\begin{remark}
The equilibrium relation:
\[
v^2\kappa
=
g\sin\phi
\]
is fundamentally a runtime-state coupling law.
\end{remark}

% ------------------------------------------------------------

\begin{remark}
It connects:
\begin{itemize}
\item intrinsic geometry,
\item cant evolution,
\item runtime vehicle states,
\item and operational admissibility.
\end{itemize}
\end{remark}

% ------------------------------------------------------------

\begin{remark}
Within AIM, dynamic balance therefore belongs naturally to runtime-state
evolution rather than to isolated geometric design formulas.
\end{remark}

% ============================================================
\section{Dynamic Balance and Realization}
% ============================================================

\begin{remark}
Operational balance depends on realized rather than purely analytical
geometry.
\end{remark}

% ------------------------------------------------------------

\begin{remark}
Thus:
\[
a_{\mathrm{eff,real}}(s)
=
v^2\kappa_{\mathrm{real}}(s)
-
g\sin\phi_{\mathrm{real}}(s)
\]
is operationally more relevant than its analytical target counterpart.
\end{remark}

% ------------------------------------------------------------

\begin{remark}
Realization therefore directly influences:
\begin{itemize}
\item comfort,
\item admissibility,
\item operational smoothness,
\item and dynamic balance.
\end{itemize}
\end{remark}

% ============================================================
\section{Dynamic Balance and Curvature Space}
% ============================================================

\begin{remark}
Operational balance emerges fundamentally from curvature evolution.
\end{remark}

% ------------------------------------------------------------

\begin{remark}
Different transition families therefore generate different operational
balance trajectories even if their position-space geometry appears
similar.
\end{remark}

% ------------------------------------------------------------

\begin{remark}
Curvature space is therefore simultaneously:
\begin{itemize}
\item a geometric domain,
\item a runtime domain,
\item and a dynamic balance domain.
\end{itemize}
\end{remark}

% ============================================================
\section{Towards Runtime-State Design}
% ============================================================

\begin{remark}
The dynamic-balance interpretation transforms railway alignment from:
\begin{quote}
curve design
\end{quote}
toward:
\begin{quote}
design of admissible runtime operational states.
\end{quote}
\end{remark}

% ------------------------------------------------------------

\begin{remark}
This prepares the transition toward:
\begin{itemize}
\item Vienna reinterpretation,
\item realization-aware dynamics,
\item Schwerpunkt-Trassierung,
\item and operational state optimization.
\end{itemize}
\end{remark}

% ============================================================
\section{Canonical Interpretation}
% ============================================================

\begin{proposition}
Within AIM, the equilibrium relation is interpreted as a runtime-state
coupling law.
\end{proposition}

% ------------------------------------------------------------

\begin{proposition}
Operational railway balance emerges from coupled runtime interaction
between:
\begin{itemize}
\item curvature,
\item cant,
\item realization,
\item speed,
\item and vehicle-state response.
\end{itemize}
\end{proposition}

% ------------------------------------------------------------

\begin{proposition}
Dynamic admissibility is fundamentally a runtime-state property rather
than a purely geometric property.
\end{proposition}

% ============================================================
\section{Connection to AIM}
% ============================================================

\begin{summary}
Dynamic balance extends railway alignment interpretation from geometric
equilibrium formulas toward runtime-state coupling.

Within AIM:
\begin{itemize}
\item equilibrium becomes a runtime-state relation,
\item cant deficiency becomes an operational imbalance state,
\item admissibility becomes envelope-based,
\item operational balance evolves dynamically,
\item and runtime-state trajectories become primary engineering objects.
\end{itemize}

This establishes the conceptual bridge between:
\begin{itemize}
\item curvature evolution,
\item runtime dynamics,
\item operational admissibility,
\item realization-aware behaviour,
\item and operational state optimization.
\end{itemize}
\end{summary}
